\documentclass[11pt,reqno]{amsart}
\usepackage[T1]{fontenc}
\usepackage{lmodern,mathrsfs,amsmath,amssymb,amsthm,mathtools,microtype,booktabs}
\usepackage[margin=1.08in]{geometry}
\usepackage{xr-hyper}
\usepackage[hidelinks]{hyperref}
\numberwithin{equation}{section}
\newtheorem{theorem}{Theorem}[section]
\newtheorem{lemma}[theorem]{Lemma}
\newtheorem{proposition}[theorem]{Proposition}
\newtheorem{corollary}[theorem]{Corollary}
\theoremstyle{definition}\newtheorem{definition}[theorem]{Definition}
\theoremstyle{remark}\newtheorem{remark}[theorem]{Remark}
\newcommand{\R}{\mathbb R}
\newcommand{\C}{\mathbb C}
\newcommand{\Pj}{\mathbb P}
\newcommand{\PP}{\mathcal P}
\newcommand{\Sym}{\operatorname{Sym}}
\newcommand{\Gr}{\operatorname{Gr}}
\newcommand{\diag}{\operatorname{diag}}
\newcommand{\rank}{\operatorname{rank}}
\newcommand{\ran}{\operatorname{ran}}
\newcommand{\Span}{\operatorname{span}}
\newcommand{\dist}{\operatorname{dist}}
\newcommand{\codim}{\operatorname{codim}}
\newcommand{\tr}{\operatorname{tr}}
\newcommand{\LL}{\mathcal L}
\newcommand{\eps}{\varepsilon}
\DeclareMathOperator*{\argmin}{argmin}
\allowdisplaybreaks[2]
\setlength{\emergencystretch}{3em}
\title[Determinantal flags and multiplication failure]{Determinantal flags and primary boundaries\\in multiplication failure}
\author{Qian Qi}
\date{September 22, 2026}
\subjclass[2020]{14M12, 14H50, 13C40, 14B05, 14N15}
\keywords{Multiplication of sections, conductor, finite algebras, primary decomposition, nonreduced schemes}
\hypersetup{pdftitle={Determinantal flags and primary boundaries in multiplication failure},pdfauthor={Qian Qi}}
